% Background and related work.
%
% Carried over from the previous version and updated in two places: the gap
% argument survives unchanged, but the paper's claim is now a finding rather
% than the suite, so §2.4 and §2.5 situate the two findings against the
% practices they amend. Every entry cited here is in refs.bib.

\section{Background and related work}
\label{sec:related}

\subsection{World models and the RSSM}
\label{sec:related:worldmodels}

A world model is an agent's internal, learned model of environment dynamics.
In the formulation of \citet{ha2018worldmodels} it comprises a visual encoder,
a transition model $M$, and a controller trained inside the model's
imagination. The recurrent state-space model \citep{hafner2019planet} gives $M$
the form used throughout this paper --- a deterministic recurrent state carried
by a GRU, and a stochastic latent sampled from it --- and the Dreamer line of
work \citep{hafner2020dreamer, hafner2021dreamerv2, hafner2023dreamerv3} builds
agents on top of it.

The relevance of $M$ to continual learning is specific: a transition model that
has forgotten a previous environment's dynamics produces incoherent imagined
rollouts for that environment \emph{regardless of how good the policy is}, and
any controller subsequently trained in that imagination inherits the corruption.
\citet{ha2018worldmodels} name catastrophic forgetting as an open problem for
world models and leave it to future work. Later systems mitigate it
incidentally: DreamerV3 \citep{hafner2023dreamerv3} trains from large replay
buffers, which mixes experience across whatever the agent has seen, but this is
a property of the data pipeline rather than a treatment of the transition
component, and it says nothing about what happens to $M$ when the mixing is
removed.

\subsection{Catastrophic forgetting and the three families of mitigation}
\label{sec:related:forgetting}

Catastrophic forgetting \citep{mccloskey1989catastrophic, ratcliff1990connectionist,
french1999catastrophic} follows from the distributed, overlapping nature of
neural representations: gradient steps that reduce loss on a new task move
shared parameters away from optima for previous ones. The stability--plasticity
dilemma \citep{vandeven2024continual} states the tension that any mitigation has
to price --- what reduces forgetting tends to reduce the capacity to learn what
comes next.

Mitigations fall into three families, and we evaluate one representative of
each. \emph{Regularization} methods add a penalty that discourages movement in
parameters deemed important for previous tasks, weighted by an importance
estimate: a diagonal Fisher in EWC \citep{kirkpatrick2017ewc}, a path integral
in synaptic intelligence \citep{zenke2017si}. \emph{Architectural} methods
allocate separate capacity per task, whether by growing it
\citep{rusu2016progressive} or by carving it out of a fixed network
\citep{mallya2018packnet}. \emph{Replay} methods retain and re-sample data from
previous tasks.

All three were developed and are predominantly evaluated in classification, or
on end-to-end reinforcement learning agents. That matters here for a reason that
is not merely about coverage: an importance estimate is defined relative to a
loss, and therefore has support only where that loss has gradients. When such a
method is applied to a world model, which loss it is attached to decides which
components it can protect at all --- a point Section~\ref{sec:results:encoder}
turns into a measurement rather than an argument.

\subsection{What existing continual learning benchmarks measure}
\label{sec:related:benchmarks}

The classification suites --- Split-MNIST, Permuted-MNIST, Split-CIFAR --- measure
forgetting in discriminative models with no temporal structure, and do not
transfer to a generative model of dynamics.

Within reinforcement learning, Continual World
\citep{wolczyk2021continualworld} is the standard benchmark: sequences of
robotic manipulation tasks with actor-critic agents, reporting forgetting,
forward transfer and performance at the level of the \emph{policy}. Closest to
our setting, \citet{kessler2023effectiveness} study world models in continual
reinforcement learning directly, evaluating DreamerV2 as an integrated system
and finding that reservoir sampling in the replay buffer is the most effective
mitigation, again measured by policy-level return, forgetting and transfer.

Both measure the agent. Neither isolates $M$, and neither defines a metric for
the quality of imagined dynamics as such. We are not aware of a benchmark that
does, and that gap is what the suite in Section~\ref{sec:method} is built to
fill. The distinction is not bookkeeping: a system-level result cannot say
whether a mitigation preserved the dynamics or merely preserved a policy that
was robust to their degradation, and those two call for different fixes.

Component-level measurement also sharpens a comparison we can now make.
\citet{kessler2023effectiveness} report that L2 regularization on the full
DreamerV2 system performs poorly. We find that EWC restricted to the transition
component preserves it almost exactly (Section~\ref{sec:results:encoder}) while
leaving everything else unprotected --- and, because its Fisher is identically
zero outside the transition loss, this is structural rather than a matter of
tuning. Read together, the two results suggest that the disappointing behaviour
of regularization in world models is less about the penalty than about the
mismatch between where its importance signal has support and where the
forgetting actually happens.

\subsection{Task distance as an experimental factor}
\label{sec:related:distance}

Continual learning results are routinely reported against a notion of how
related the tasks in a sequence are, and benchmarks are built to span that
notion: sequences are assembled from tasks judged similar or dissimilar, and the
resulting ordering is then treated as an experimental factor against which
forgetting is expected to grow. Our own design does exactly this, with three
levels of dynamic distance per family
(Section~\ref{sec:method:families}), and the previous version of this work
described them as controlled.

What is usually missing --- and was missing from our own design --- is a
\emph{measured} quantity to which the ordering is answerable. Where the
environments differ in explicit physical parameters, differencing them is
available and exact (Equation~\ref{eq:dparam}); where they differ in body or
layout, there is nothing to difference, and the ordering rests on judgement.
Section~\ref{sec:results:axis} reports what happened when we checked ours
against both the forgetting it produced and a measured distance between the
environments' transition models, and Section~\ref{sec:discussion:axis} states
the recommendation that follows. The recommendation is aimed at this practice
rather than at any particular benchmark.

\subsection{Where forgetting is measured inside a model}
\label{sec:related:locus}

A metric for a multi-component model has to hold something fixed in order to
attribute a change to anything, and what it holds fixed bounds what it can
detect. Metrics defined over a representation --- ours included, since PF and RD
score the transition component in a latent basis encoded once
(Section~\ref{sec:method:metrics}) --- are by construction unable to observe
drift in the mapping that produced that basis. This is the price of component
isolation, not a defect of a particular estimator, and the same trade appears
wherever forgetting is attributed to a part of a network rather than to the
whole.

The consequence is that a component-level benchmark owes its readers both
scales. Reporting only the isolated metric can certify that a method preserved
the component it was pointed at while the model as a whole became unusable for
the old task. Section~\ref{sec:results:encoder} is that situation, measured.
